\section*{Author Section 4: improved PRG for constant-depth circuits}
We give the formal statement of the derandomized multi-switching lemma below. This is the full version of Lemma 3 in the introduction.
\paragraph{Lemma 9.} Let $t,w,k,n,m\ge1$ be integers. Let $p,\epsilon>0$ be reals. Let $\mathcal F=\{F_1,\ldots,F_m\}$ be a list of size-$m$ $k$-DNFs on inputs $\{0,1\}^n$. Let $(\boldsymbol\Lambda,\boldsymbol x)$ be a joint random variable satisfying the following:
\begin{itemize}\item $\boldsymbol\Lambda$ is a $(t+k)$-wise $p$-bounded subset of $[n]$.
\item Conditioning on $\boldsymbol\Lambda$, $\boldsymbol x$ is an $n$-bit random string that $\epsilon$-fools size-$(m^2)$ CNF.\end{itemize}
Then we have
\begin{multline*}\Pr_{T\sim\boldsymbol\Lambda,x\sim\boldsymbol x}[\mathcal F|_{\Lambda[x,\star]}\text{ does not have }w\text{-partial depth-}t\text{ DT}]\\\le4m^{t/w}(24pk)^t+(24m)^{t+k}\epsilon.\end{multline*}
We defer the proof of Lemma 9 to Section 6. Assuming Lemma 9, we prove Theorem 1. We prove a slightly stronger form of Theorem 1, stated below.
\paragraph{Theorem 10.} For every $d\ge2$ the following is true. For every $m,n\in\mathbb N$ such that $m\ge n$, $k\in\mathbb N$ and every $\epsilon>0$, there is an $\epsilon$-error, $O((\log^2(m)+k\log^{d-2}(m))\log(m/\epsilon)\log\log(m))$-seed PRG for $\mathrm{AC}^0_d[m]$ circuits with bottom fan-in bounded by $k$.
Theorem 10 shows that, when the bottom fan-in of the $\mathrm{AC}^0_d$ circuits is smaller than $o(\log(m))$, we can hope for shorter seed length through our construction. Note that we can always interpret a depth-$d$ $\mathrm{AC}^0$ circuit with unbounded bottom fan-in as a depth-$(d+1)$ $\mathrm{AC}^0$ circuit with bottom fan-in being 1. Then, Theorem 1 follows from Theorem 10 easily.
\begin{proof}We use induction on the depth $d$. The case for $d=2$ follows from Theorem 8. Assuming this is true for depth $d-1\ge2$, we prove it for the case of $d$. Let $w=40k$ and $t=80\log(m/\epsilon)$. We prepare the following pseudorandom primitives.
\begin{itemize}\item First, let $H:[n]\to[w]$ be a $2t$-wise independent hash function, samplable using $O(\log(n)\log(m/\epsilon))$ bits. In the following, we will use $H_i$ to denote $H^{-1}(i)$. We remark that $H_i$ can be equivalently expressed as an $n$-bit string. Namely, $(H_i)_j=1$ if and only if $j\in H_i$.
\item Second, let $\epsilon'=\epsilon/(w2^{t+1})$. Sample $X_1,\ldots,X_w\in\{0,1\}^n$, each being an independent string that $\epsilon'$-fools $\mathrm{AC}^0_{d-1}$-circuits of size $2m^2$ and bottom-width $\log m$. When $d-1\ge3$, $X_i$ is samplable using $O(\log^{d-2}(m)\log(m/\epsilon)\log\log(m))$ bits by the induction hypothesis. For the case $d-1=2$, $X_i$ is samplable using $O(\log(m)\log(m/\epsilon)\log\log(m))$ bits by Theorem 8.
\item Lastly, let $Y\in\{0,1\}^n$ be a random string that $\epsilon/(24m)^{2t}$-fools size-$m$ CNF, samplable using $O(\log^2m\log(m/\epsilon)\log\log m)$ bits (by Theorem 8).\end{itemize}
The seed for our generator is the concatenation of the seeds used to sample all the primitives above. We compute the output of our generator as
\[Y\oplus(X_1\land H_1)\oplus(X_2\land H_2)\oplus\cdots\oplus(X_w\land H_w).\tag{1}\]
Here, $\land$ and $\oplus$ denote bit-wise AND and XOR respectively. The seed length is bounded by
\begin{align*}&O(\log(n)\log(m/\epsilon)+w\log^{d-2}(m)\log(m2^t/\epsilon)\log\log(m)\\&\qquad+\log^2m\log(m/\epsilon)\log\log m)\\&\le O((\log^2(m)+k\log^{d-2}(m))\log(m/\epsilon)\log\log(m)).\end{align*}
We argue the correctness by a hybrid argument. Fix $C$ to be an $\mathrm{AC}^0_d$-circuit that we wish to fool. Let $U_1,\ldots,U_w$ denote $w$ independent uniformly random strings from $\{0,1\}^n$. For every $i\in\{0,1,\ldots,w\}$, we define the $i$-th hybrid distribution as
\begin{multline*}D_i:=Y\oplus(U_1\land H_1)\oplus\cdots\oplus(U_i\land H_i)\\\oplus(X_{i+1}\land H_{i+1})\oplus\cdots\oplus(X_w\land H_w).\tag{2}\end{multline*}
We observe that $D_0$ is the output distribution of our PRG, while $D_w$ is a uniformly random string from $\{0,1\}^n$. Hence, it suffices to show that
\[|\mathbb E_{x\sim D_0}[C(x)]-\mathbb E_{x\sim D_w}[C(x)]|\le\epsilon.\tag{3}\]
To show (3), it suffices to show for every $i\in\{1,\ldots,w\}$ that
\[|\mathbb E_{x\sim D_{i-1}}[C(x)]-\mathbb E_{x\sim D_i}[C(x)]|\le\epsilon/w.\tag{4}\]
In the following, we prove (4). We observe that $H_i$ is $2t$-wise $1/w$-bounded. Conditioning on an instantiation of $H$, we have that $Z_i:=Y\oplus\sum_{j<i}(U_j\land H_j)\oplus\sum_{j>i}(X_j\land H_j)$ is an $\epsilon/(24m)^{2t}$-pseudorandom string for size-$(m^2)$ CNF, because $Y$ is. Let $\mathcal F$ be the family of all next-to-bottom layer sub-circuits of $C$. Denote by $E$ the event ``$\mathcal F|_{H_i[Z_i,\star]}$ does not have $\log(m)$-partial depth-$t$ DT.'' Then it follows from Lemma 9 that
\begin{align*}&\Pr_{H,Y,U_1,\ldots,U_{i-1},X_{i+1},\ldots,X_w}[E]\\&\qquad\le4m^{t/\log(m)}(24k/w)^t+\frac{\epsilon(24m)^{t+\log(m)}}{(24m)^{2t}}\le\frac\epsilon{2w}.\end{align*}
Conditioning on $\neg E,H,Y,U_1,\ldots,U_{i-1},X_{i+1},\ldots,X_w$ and calculating $Z_i$ as defined above, one can then write $C|_{H_i[Z_i,\star]}$ as a depth-$t$ decision tree $T$ where each leaf of $T$ is labelled by an $\mathrm{AC}^0_{d-1}$-circuit of size $m^2$ and bottom fan-in $\log(m)$. Let $\{\ell_1,\ldots,\ell_{2^t}\}$ enumerate the leaves of the decision tree. Each $\ell_j$ is associated with a size-$m^2$ depth-$(d-1)$ circuit, which is also denoted by $\ell_j:\{0,1\}^n\to\{0,1\}$ for brevity. Then one can write $C|_{(H_i)[Z_i,\star]}$ as
\[C|_{H_i[Z_i,\star]}(y)=\sum_{j=1}^{2^t}\ell_j(y)\mathbf1_{\{T(y)\text{ reaches leaf }\ell_j\}}.\]
Let's fix an index $j\in[2^t]$ for now. Note that $\ell_j(y)\mathbf1_{\{\text{reach }\ell_j\text{ on }y\}}$ is itself a depth-$(d-1)$ circuit of size at most $2m^2$. By the construction of $X_i$ we know that
\begin{multline*}\mathbb E_{X_i}[\ell_j(X_i+Z_i)\mathbf1_{\{T\text{ reaches }\ell_j\text{ on }X_i+Z_i\}}]-\\\mathbb E_{U_i}[\ell_j(U_i+Z_i)\mathbf1_{\{T\text{ reaches }\ell_j\text{ on }U_i+Z_i\}}]\le\frac\epsilon{2^{t+1}w}.\end{multline*}
Taking a summation over all leaves $j$, one gets
\begin{multline*}\mathbb E_{X_i}[C|_{H_i[Z_i,\star]}(X_i+Z_i)]-\mathbb E_{U_i}[C|_{H_i[Z_i,\star]}(U_i+Z_i)]]\\\le\epsilon/(2w).\end{multline*}
Finally, one has
\[|\mathbb E_{x\sim D_{i-1}}[C(x)]-\mathbb E_{x\sim D_i}[C(x)]|\le\Pr[\neg E]\frac\epsilon{2w}+\Pr[E]\le\frac\epsilon w,\]
proving (4).\end{proof}
Given Theorem 10, we prove Theorem 1 by tuning parameters.
\paragraph{Proof of Theorem 1.} For every $d\ge3$ and every $\mathrm{AC}^0_d[m]$ circuit with unbounded bottom fan-in, we can interpret it as a depth-$(d+1)$ circuit with bottom fan-in being 1. Applying Theorem 10 in this case gives a PRG with seed length $O(\log^{d-1}(m)\log(m/\epsilon)\log\log(m))$, as desired. For the case of $d=2$, we use Theorem 8 directly. This completes the proof.\qed
As a final remark, suppose we could have $O(\log(m)\log(m/\epsilon))$-seed PRG for CNFs (namely, if we can shave the $\log\log m$ factor in Theorem 8). Then our construction implies PRG for $\mathrm{AC}^0_3$ with seed length $O(\log^2(m)\log(m/\epsilon))$, and further implies PRG for $\mathrm{AC}^0_d$ with seed length $O(\log^{d-1}(m)\log(m/\epsilon))$, matching the lower bound barrier [14,13].
